% ─────────────────────────────────────────────────────────────────────────────
\section{Feature Engineering}
\label{sec:features}
% ─────────────────────────────────────────────────────────────────────────────

A central challenge in applying classical ML to variable-length time series is
the extraction of a fixed-size feature vector that captures diagnostically
relevant information while discarding operational noise. This section describes
the hierarchy of representations developed for NGAFID, from a simple global
statistics baseline to segmented and cross-physical descriptors.

\subsection{Representation B: Global Per-Sensor Statistics}
\label{subsec:reprB}

The foundational representation compresses each flight's
$\mathbf{X} \in \mathbb{R}^{2048 \times 23}$ into a fixed vector by computing
eight aggregate statistics per sensor channel:

\begin{equation}
    \phi(\mathbf{x}_j) = \bigl[\mu_j,\ \sigma_j,\ x_j^{\min},\ x_j^{\max},\
        r_j,\ \text{skew}_j,\ \text{kurt}_j,\ \hat{\beta}_j\bigr]
    \label{eq:reprB_stats}
\end{equation}

where $\mathbf{x}_j \in \mathbb{R}^{2048}$ is the time series for sensor $j$,
$\mu_j$ is the mean, $\sigma_j$ the standard deviation, $r_j = x_j^{\max} -
x_j^{\min}$ the range, and $\hat{\beta}_j$ the ordinary least-squares slope of
a linear trend fit to the time series (computed via the closed-form formula
$\hat{\beta} = \frac{\sum_t (t - \bar{t})(x_t - \mu)}{\sum_t (t - \bar{t})^2}$
for efficiency).

Concatenating these eight statistics over all 23 channels yields:
\begin{equation}
    \mathbf{f}^{(B)} \in \mathbb{R}^{184}, \quad
    184 = 23 \times 8
    \label{eq:reprB_dim}
\end{equation}

\paragraph{Missing values.} Raw features contain 1{,}164 NaN values (across the
full 11{,}446-flight dataset), arising from sensor drop-outs on individual
flights. These are imputed with the column-wise median of the training fold,
reducing NaN count to zero before model fitting.

\paragraph{Build time.} Computing Repr.~B for 11{,}446 flights takes
approximately 185\,s on a standard laptop. The linear slope is computed using a
fast closed-form expression (approximately $5\times$ faster than
\texttt{numpy.polyfit}).

\paragraph{PCA analysis.} When Repr.~B is projected into principal component
space, \textbf{51 components} are needed to explain 95\% of the variance. This
indicates a high-dimensional structure with no single dominant direction, and
confirms that fault signatures---concentrated in PCs 12--18 at $<$0.1\% of
total variance---are effectively buried by the operational variability captured
in the top components.

\subsection{Representation B\textsubscript{seg}: Temporal Segmentation}
\label{subsec:reprBseg}

Repr.~B discards all temporal structure: a mean of 3{,}500\,RPM is the same
whether the engine stayed at that value throughout the flight or ramped up and
then down. To reintroduce \textit{temporal ordering} while remaining within the
classical ML framework, each flight is divided into $N$ equal-length segments
and Repr.~B statistics are computed independently on each segment.

Formally, for a flight of 2{,}048 steps divided into $N$ segments of $L = 2048/N$
steps each, the segmented feature vector is:
\begin{equation}
    \mathbf{f}^{(B_{\text{seg}})} = \bigl[\phi(\mathbf{x}_j^{(1)}),\ \ldots,\
        \phi(\mathbf{x}_j^{(N)})\bigr]_{j=1}^{23}
        \in \mathbb{R}^{23 \times 8 \times N}
\end{equation}

Table~\ref{tab:reprBseg} summarizes the feature counts and temporal granularity
for each segmentation level.

\begin{table}[H]
    \centering
    \caption{Repr.~B\textsubscript{seg} configurations.}
    \label{tab:reprBseg}
    \begin{tabular}{cccc}
        \toprule
        $N_{\text{seg}}$ & \textbf{Features} & \textbf{Temporal resolution} & \textbf{Segment length} \\
        \midrule
        1  &    184 & None (global) & 2{,}048 steps \\
        3  &    552 & Thirds of flight & 682 steps \\
        5  &    920 & Fifths of flight & 409 steps \\
        10 & 1{,}840 & Tenths of flight & 204 steps \\
        \bottomrule
    \end{tabular}
\end{table}

The hypothesis is that the \textit{evolution} of sensor statistics across
flight thirds or fifths carries diagnostic information beyond the global
distribution.

\subsection{Representation F\textsubscript{++}: Cross-Physical Features}
\label{subsec:reprFpp}

Beyond per-sensor statistics, physically meaningful relationships between
sensors can provide additional discriminative power. Repr.~F\textsubscript{++}
extends Repr.~B with a set of \textit{inter-cylinder features}:

\begin{itemize}
    \item \textbf{Spread:} Maximum minus minimum across the four CHT (and four
          EGT) channels at each time step, then summarized with the same eight
          statistics. A large CHT spread indicates uneven combustion or a
          failing cylinder---a fault signature invisible to per-channel averages.
    \item \textbf{Coefficient of variation:} $\text{CV} = \sigma / \mu$ for
          CHT and EGT cross-cylinder variability.
    \item \textbf{Cross-cylinder ratios:} Pairwise ratios (e.g.,
          CHT1/CHT2, EGT1/EGT2) under statistical summaries.
    \item \textbf{Extended statistics on extremes:} Slope and range computed on
          the per-window max and min envelopes, capturing transient events missed
          by global statistics.
\end{itemize}

This brings the total feature count to \textbf{224 features}. For the FC task,
XGBoost gain-based feature selection further reduces this to \textbf{60 features}
(Repr.~F), eliminating low-information channels while preserving the most
discriminative ones. The top 5 features by XGBoost information gain are:
\texttt{E1 FFlow\_min}, \texttt{OAT\_min}, \texttt{amp2\_mean},
\texttt{amp2\_min}, and \texttt{volt2\_min}.

Table~\ref{tab:repr_summary} summarizes all representations and their key
properties.

\begin{table}[H]
    \centering
    \caption{Summary of feature representations.}
    \label{tab:repr_summary}
    \begin{tabular}{llcp{5.5cm}}
        \toprule
        \textbf{Repr.} & \textbf{Method} & \textbf{Dim.} & \textbf{What it captures} \\
        \midrule
        B                & Global stats (8 per sensor)  & 184     & Magnitude, spread, trend of full flight \\
        B\textsubscript{seg} & B applied to $N$ segments & 552--1840 & Temporal evolution of global stats \\
        F\textsubscript{++} & B + cross-cylinder features & 224   & Inter-cylinder asymmetry and ratios \\
        F               & F\textsubscript{++} + XGB gain selection & 60 & Most discriminative 60 features \\
        \bottomrule
    \end{tabular}
\end{table}
